% Einheit 1 von 3 – Summe mal Summe (bank/binomische-formeln/e1.jsonl, zusammenbau v0.1)
%% TODO Zweigzeile Teil 1: was hier gelernt wird (Satz in Schülersprache aus dem Einheitstitel) – Einheit 1
\einheitenkopf[e1]{Einheit 1 von 3 · Summe mal Summe}
\zweigzeile{ab Kl. 8, je nach Buch bis Kl. 9 (am Gymnasium ab Kl. 7) · keine P10-Aufgabe}
\setcounter{aufgabe}{9}

%% TODO Titel: Ich-kann-Satz fehlt in der Bank (Platzhalter: Kettenname) – Nr. 10
%% TODO Anweisung: ein Satz über der ersten Teilaufgabe fehlt in der Bank – Nr. 10
\begin{aufgabe}{Summe mal Summe}
\begin{teile}
\teil Verbinde jedes Glied der ersten Klammer mit jedem Glied der zweiten durch einen Pfeil und schreib die Produkte hin, ohne zu rechnen: $(x + 4) \cdot (y + 6)$ \leerfeld
\end{teile}
%% TODO Darstellung für schwach fehlt in der Bank (binomische-formeln-e1-k1-s1-v1; Abschnitt „Für schwache Schüler“)
\swz{Fasse zusammen: $3x + 4y + 2x + y$}{}
%% TODO Darstellung für schwach fehlt in der Bank (binomische-formeln-e1-k1-s1-v2; Abschnitt „Für schwache Schüler“)
\swz{Fasse zusammen: $2a + 5b + 4a + 3b$}{}
%% TODO Darstellung für schwach fehlt in der Bank (binomische-formeln-e1-k1-s1-v3; Abschnitt „Für schwache Schüler“)
\swz{Fasse zusammen: $x + 6y + 7x + 2y$}{}
%% TODO Darstellung für schwach fehlt in der Bank (binomische-formeln-e1-k1-s1-v4; Abschnitt „Für schwache Schüler“)
\swz{Fasse zusammen: $5y + 2x + 3y + 4x$}{}
%% TODO Darstellung für schwach fehlt in der Bank (binomische-formeln-e1-k1-s1-v5; Abschnitt „Für schwache Schüler“)
\swz{Fasse zusammen: $7a + 2b + a + 9b$}{}
\swfrage{Was bleibt gleich, was ändert sich?}
%% TODO Darstellung für schwach fehlt in der Bank (binomische-formeln-e1-k1-s2-v1; Abschnitt „Für schwache Schüler“)
\swz{Fasse zusammen: $6x - 2y - x + 5y$}{}
%% TODO Darstellung für schwach fehlt in der Bank (binomische-formeln-e1-k1-s3-v1; Abschnitt „Für schwache Schüler“)
\swz{Fasse zusammen: $2x^2 + 3xy + x^2 + 4x - xy$}{}
%% TODO Darstellung für schwach fehlt in der Bank (binomische-formeln-e1-k1-s4-v1; Abschnitt „Für schwache Schüler“)
\swz{$\text{Multipliziere aus: $4 \cdot (2x + 3y)$}$}{}
%% TODO Darstellung für schwach fehlt in der Bank (binomische-formeln-e1-k1-s5-v1; Abschnitt „Für schwache Schüler“)
\swz{$\text{Schreib die vier Produkte hin: $(x + 4) \cdot (x + 7)$}$}{}
%% TODO Darstellung für schwach fehlt in der Bank (binomische-formeln-e1-k1-s6-v1; Abschnitt „Für schwache Schüler“)
\swz{$\text{Multipliziere aus: $(x + 2) \cdot (x + 8)$}$}{}
%% TODO Darstellung für schwach fehlt in der Bank (binomische-formeln-e1-k1-s7-v1; Abschnitt „Für schwache Schüler“)
\swz{$\text{Multipliziere aus: $(x + 6) \cdot (x - 2)$}$}{}
%% TODO Darstellung für schwach fehlt in der Bank (binomische-formeln-e1-k1-s8-v1; Abschnitt „Für schwache Schüler“)
\swz{$\text{Multipliziere aus: $(x - 6) \cdot (x - 1)$}$}{}
%% TODO Darstellung für schwach fehlt in der Bank (binomische-formeln-e1-k1-s9-v1; Abschnitt „Für schwache Schüler“)
\swz{$\text{Multipliziere aus: $(2x + 1) \cdot (x + 4)$}$}{}
%% TODO Darstellung für schwach fehlt in der Bank (binomische-formeln-e1-k1-s10-v1; Abschnitt „Für schwache Schüler“)
\swz{$\text{Multipliziere aus: $(a + 2b) \cdot (3a + b)$}$}{}
\swz[3]{$\text{Multipliziere aus und fasse zusammen: $(3x - 4) \cdot (2x + 5)$}$}{}
\end{aufgabe}

%% TODO Titel: Ich-kann-Satz fehlt in der Bank (Platzhalter: Kettenname) – Nr. 11
%% TODO Anweisung: ein Satz über der ersten Teilaufgabe fehlt in der Bank – Nr. 11
\begin{aufgabe}{Termwert mit zwei Variablen berechnen (auch negativ)}
%% TODO Darstellung für schwach fehlt in der Bank (binomische-formeln-e1-k2-s1-v1; Abschnitt „Für schwache Schüler“)
\swz{$3x + 2y$ für $x = 4$ und $y = -5$ – Termwert?}{}
\end{aufgabe}

%% TODO Titel: Ich-kann-Satz fehlt in der Bank (Platzhalter: Kettenname) – Nr. 12
%% TODO Anweisung: ein Satz über der ersten Teilaufgabe fehlt in der Bank – Nr. 12
\begin{aufgabe}{Summe mal Summe – Fehler finden, Begründen}
\swz[3]{Finde den Fehler und rechne richtig: \rechnung{(x + 6) \cdot (x + 4) &= x^2 + 24}}{}
\swfrage{Warum entstehen bei $(x + 2) \cdot (y + 5)$ vier Produkte? Begründe am Flächenbild eines Rechtecks.}
\end{aufgabe}

\uebersichtskasten{Klammer mal Klammer: jedes Glied der ersten Klammer mal jedes Glied der zweiten (vier Produkte), dann zusammenfassen. \\ (x + 3)·(x + 5) = x² + 5x + 3x + 15 = x² + 8x + 15 \\ Vorzeichen gehören zum Glied: (x − 4)·(x + 2) = x² + 2x − 4x − 8 = x² − 2x − 8 \\ Zwei Variablen: nur gleiche Variablen mit gleicher Hochzahl zusammenfassen; xy und yx sind dasselbe. \quad (a + b)·(c + d) = ac + ad + bc + bd}
